%=====================================================================
\chapter{Deep Architectures}\label{ch:architectures}
%=====================================================================
\locator{4}{Part IV --- Deep Learning and Modern AI}

\begin{outcomes}
By the end of this chapter you will be able to:
\begin{tight}
  \item \textbf{Explain} what convolution does and why it suits grid-structured data.
  \item \textbf{Describe} how recurrent networks carry state across a sequence, and why LSTMs were needed.
  \item \textbf{State} the core idea of attention and why it displaced recurrence.
  \item \textbf{Use} an autoencoder for compression and anomaly detection.
  \item \textbf{Apply} transfer learning and justify it on a small dataset.
\end{tight}
\end{outcomes}

\begin{prereq}
\chref{neuralnets} --- neurons, layers, activations, backpropagation. Everything
here is that machinery with the connections arranged differently.
\end{prereq}

\begin{keyterms}
convolution $\bullet$ kernel / filter $\bullet$ feature map $\bullet$ pooling
$\bullet$ CNN $\bullet$ recurrent network $\bullet$ hidden state $\bullet$ LSTM
$\bullet$ GRU $\bullet$ attention $\bullet$ self-attention $\bullet$ transformer
$\bullet$ embedding $\bullet$ autoencoder $\bullet$ latent space $\bullet$ transfer
learning $\bullet$ fine-tuning
\end{keyterms}

\section{Architecture Is Inductive Bias}

\begin{conceptbox}[The Idea Behind This Chapter]
A fully connected network treats every input the same way and must learn all
structure from data. An \emph{architecture} builds a structural assumption into the
wiring: convolution assumes nearby pixels are related; recurrence assumes order
matters; attention assumes any element may relate to any other. Choosing an
architecture is choosing which assumption to give the model for free --- which is
why the right one needs far less data.
\end{conceptbox}

\section{Convolutional Networks}

\begin{definitionbox}[Convolution]
Slide a small matrix of weights --- a \term{kernel} or \term{filter} --- across the
input, computing a dot product at each position. The output is a \term{feature map}
showing where in the input that pattern occurs. The same weights are reused at
every position, which is called \term{weight sharing}.
\end{definitionbox}

\dsfig{%
\begin{tikzpicture}[font=\scriptsize,
  c/.style={draw=gray!55, minimum size=0.44cm, inner sep=0pt, font=\tiny},
  k/.style={draw=accent, fill=accent!16, minimum size=0.44cm, inner sep=0pt, font=\tiny},
  o/.style={draw=greenhl, fill=greenhl!14, minimum size=0.44cm, inner sep=0pt, font=\tiny}]

% input grid 5x5
\node[font=\scriptsize\bfseries, text=primary] at (1.1,2.7) {input};
\foreach \r in {0,...,4}{\foreach \cc in {0,...,4}{
  \pgfmathsetmacro{\v}{mod(\r*3+\cc*2,5)}
  \pgfmathtruncatemacro{\vi}{\v}
  \node[c] at (\cc*0.45, 2.2-\r*0.45) {\vi};}}
% highlight receptive field
\draw[accent, line width=1.3pt] (-0.24,2.44) rectangle (1.11,1.09);

% kernel
\node[font=\scriptsize\bfseries, text=accent!75!black] at (3.5,2.7) {kernel $3{\times}3$};
\foreach \r in {0,...,2}{\foreach \cc in {0,...,2}{
  \node[k] at (3.0+\cc*0.45, 2.2-\r*0.45) {\ifnum\r=1 1\else 0\fi};}}
\node[font=\tiny, text=gray!55!black, align=center, text width=2.2cm] at (3.45,0.35)
     {this one detects\\horizontal edges};

% output
\node[font=\scriptsize\bfseries, text=greenhl!45!black] at (6.7,2.7) {feature map};
\foreach \r in {0,...,2}{\foreach \cc in {0,...,2}{
  \node[o] at (6.2+\cc*0.45, 2.2-\r*0.45) {};}}
\draw[greenhl, line width=1.3pt] (5.98,2.44) rectangle (6.43,1.99);

\draw[-{Stealth[length=2mm]}, gray!60] (1.35,1.75) -- (2.75,1.75);
\draw[-{Stealth[length=2mm]}, gray!60] (4.3,1.75) -- (5.9,1.75);
\node[font=\tiny, text=gray!55!black] at (2.05,2.0) {$\otimes$};
\node[font=\tiny, text=gray!55!black] at (5.1,2.0) {$=$};

\node[align=left, text width=6.2cm, font=\tiny, text=gray!55!black, anchor=north west]
     at (8.0,2.6)
     {\textbf{Why this matters:}\\[3pt]
      $\bullet$ \textbf{Weight sharing} --- the same 9 weights are used everywhere, so a
      $5{\times}5$ input needs 9 parameters, not 225.\\[3pt]
      $\bullet$ \textbf{Translation invariance} --- an edge is detected wherever it
      appears.\\[3pt]
      $\bullet$ \textbf{Locality} --- each output depends only on a small neighbourhood,
      which is exactly true for images and for sensor grids.\\[3pt]
      Stacking convolutions builds a hierarchy: edges $\rightarrow$ textures
      $\rightarrow$ parts $\rightarrow$ objects.};
\end{tikzpicture}}%
{Convolution. One small kernel slides over the whole input, so the network learns
\emph{what} to look for rather than \emph{where} --- a massive reduction in parameters
compared with a fully connected layer.}{convolution}

\begin{definitionbox}[Pooling]
Downsample a feature map by taking the maximum (or mean) over small windows. This
shrinks the representation, reduces computation and grants a little tolerance to
small shifts in position.
\end{definitionbox}

\dsfig{%
\begin{tikzpicture}[font=\scriptsize,
  st/.style={rounded corners=3pt, draw=#1, fill=#1!9, text=#1!50!black,
             text width=1.85cm, align=center, font=\tiny\bfseries, minimum height=1.5cm},
  arr/.style={-{Stealth[length=2mm]}, gray!60, line width=0.9pt}]

\node[st=secondary] (a) at (0,0)    {INPUT\\{\tiny $224{\times}224{\times}3$}};
\node[st=primary]   (b) at (2.25,0) {CONV\\+ ReLU};
\node[st=goldhl]    (c) at (4.5,0)  {POOL\\{\tiny halve size}};
\node[st=primary]   (d) at (6.75,0) {CONV\\+ ReLU};
\node[st=goldhl]    (e) at (9.0,0)  {POOL};
\node[st=purplehl]  (f) at (11.25,0){FLATTEN\\+ DENSE};
\node[st=accent]    (g) at (13.5,0) {SOFTMAX\\{\tiny $k$ classes}};

\foreach \x/\y in {a/b,b/c,c/d,d/e,e/f,f/g}{\draw[arr] (\x) -- (\y);}

\node[font=\tiny\itshape, text=gray!55!black, align=center, text width=3.4cm] at (2.3,-1.35)
     {early layers:\\edges and colours};
\node[font=\tiny\itshape, text=gray!55!black, align=center, text width=3.4cm] at (7.5,-1.35)
     {middle layers:\\textures and parts};
\node[font=\tiny\itshape, text=gray!55!black, align=center, text width=3.0cm] at (12.4,-1.35)
     {late layers:\\whole objects};
\draw[-{Stealth[length=2.4mm]}, primary!30, line width=5pt, opacity=0.4]
      (-0.9,-1.95) -- (14.4,-1.95);
\node[font=\tiny\bfseries, text=primary!60!black] at (6.7,-2.28)
     {representations become more abstract and more task-specific from left to right};
\end{tikzpicture}}%
{A typical CNN pipeline. The left-hand layers learn generic visual features that transfer
to almost any image task --- which is precisely what makes transfer learning
work.}{cnn-pipeline}

\section{Recurrent Networks}

\begin{definitionbox}[Recurrent Neural Network]
An RNN processes a sequence one element at a time, maintaining a \term{hidden
state} $h_t$ that summarises everything seen so far:
\[
h_t = f(W_x x_t + W_h h_{t-1} + b)
\]
The same weights are applied at every timestep --- weight sharing across
\emph{time} rather than across space.
\end{definitionbox}

\begin{alertbox}[Why Plain RNNs Fail]
Backpropagating through 100 timesteps multiplies 100 gradients together. If they
are slightly below 1, the product vanishes and the network cannot learn long-range
dependencies; slightly above 1 and it explodes. In practice a plain RNN forgets
anything more than roughly ten steps back --- which is useless for a sentence, a
session or a day of telemetry.
\end{alertbox}

\dsfig{%
\begin{tikzpicture}[font=\scriptsize,
  cell/.style={rounded corners=3pt, draw=primary, fill=primary!10, text=primary!55!black,
               minimum width=1.25cm, minimum height=1.0cm, font=\tiny\bfseries},
  gate/.style={circle, draw=#1, fill=#1!16, text=#1!50!black, minimum size=0.6cm,
               inner sep=0pt, font=\tiny\bfseries},
  arr/.style={-{Stealth[length=1.8mm]}, gray!65, line width=0.8pt}]

% unrolled RNN
\node[font=\scriptsize\bfseries, text=primary, anchor=west] at (-0.9,1.85) {RNN unrolled through time};
\foreach \t in {1,2,3,4}{
  \pgfmathsetmacro{\xx}{(\t-1)*2.35}
  \node[cell] (c\t) at (\xx,0.55) {$h_\t$};
  \node[font=\tiny, text=secondary!55!black] (i\t) at (\xx,-0.75) {$x_\t$};
  \node[font=\tiny, text=accent!70!black] (o\t) at (\xx,1.75) {$\hat{y}_\t$};
  \draw[arr] (i\t) -- (c\t);
  \draw[arr] (c\t) -- (o\t);}
\foreach \a/\b in {1/2,2/3,3/4}{
  \draw[arr, primary!65, line width=1pt] (c\a) -- (c\b);}
\node[font=\tiny, text=primary!60!black] at (1.17,0.9) {state};
\node[font=\tiny\itshape, text=gray!55!black, anchor=west, text width=3.4cm] at (7.6,0.55)
     {the same weights are reused at every step};

% LSTM cell
\begin{scope}[shift={(0,-3.6)}]
  \node[font=\scriptsize\bfseries, text=greenhl!45!black, anchor=west] at (-0.9,1.7)
       {LSTM cell: three gates control what is kept};
  \node[rounded corners=5pt, draw=greenhl, fill=greenhl!6, minimum width=6.6cm,
        minimum height=2.0cm] (box) at (2.9,0.35) {};
  \node[gate=accent]   (fg) at (1.0,0.75)  {$f$};
  \node[gate=secondary](ig) at (2.9,0.75)  {$i$};
  \node[gate=goldhl]   (og) at (4.8,0.75)  {$o$};
  \node[font=\tiny, text=accent!70!black, align=center]    at (1.0,-0.25) {forget\\gate};
  \node[font=\tiny, text=secondary!55!black, align=center] at (2.9,-0.25) {input\\gate};
  \node[font=\tiny, text=goldhl!45!black, align=center]    at (4.8,-0.25) {output\\gate};

  \draw[arr, greenhl!70, line width=1.2pt] (-0.5,1.15) -- (6.3,1.15);
  \node[font=\tiny, text=greenhl!45!black, anchor=south] at (6.0,1.2) {cell state};
  \node[font=\tiny\itshape, text=gray!55!black, anchor=west, text width=5.6cm] at (7.0,0.35)
       {The cell state runs straight through with only additive edits, so gradients
        survive hundreds of steps. The gates \emph{learn} what to forget, what to store
        and what to expose --- which is why an LSTM can remember something from the start
        of a session.};
\end{scope}
\end{tikzpicture}}%
{Top: an RNN unrolled, carrying state forward. Bottom: an LSTM cell, whose gated cell state
provides an uninterrupted path for gradients and solves the forgetting problem.}{rnn-lstm}

\section{Attention and Transformers}

\begin{definitionbox}[Attention]
Instead of compressing a sequence into one hidden state, attention lets every
position look directly at every other position and decide how much each one
matters. Each output is a weighted average of all inputs, where the model
\emph{learns} the weights from the content.
\end{definitionbox}

\dsfig{%
\begin{tikzpicture}[font=\scriptsize,
  tok/.style={rounded corners=2pt, draw=primary!60, fill=primary!10, text=primary!55!black,
              minimum width=1.25cm, minimum height=0.5cm, font=\tiny}]

\foreach \i/\w in {0/The, 1/server, 2/that, 3/we, 4/patched, 5/crashed}{
  \node[tok] (t\i) at (\i*1.55,0) {\w};}

% attention weights from "crashed"
\node[rounded corners=2pt, draw=accent, fill=accent!16, text=accent!70!black,
      minimum width=1.25cm, minimum height=0.5cm, font=\tiny\bfseries] at (5*1.55,0) {crashed};

\foreach \i/\wt/\op in {0/0.05/0.12, 1/0.62/0.85, 2/0.04/0.10, 3/0.03/0.08, 4/0.26/0.45}{
  \draw[-{Stealth[length=1.6mm]}, accent, opacity=\op, line width=2.0pt]
       (5*1.55,-0.3) .. controls +(0,-0.8) and +(0,-0.8) .. (\i*1.55,-0.3);
  \node[font=\tiny, text=accent!70!black, opacity=\op] at (\i*1.55,-1.5) {\wt};}

\node[font=\scriptsize\itshape, text=gray!55!black, anchor=north, align=center, text width=11cm]
     at (3.9,-2.0)
     {To interpret \emph{crashed}, the model attends most to \emph{server} (0.62) and
      \emph{patched} (0.26) --- and it learned that, rather than being told. The distance
      between \emph{server} and \emph{crashed} is irrelevant: attention reaches any
      position in one step, where an RNN would need five.};

\node[anchor=north west, align=left, text width=13.2cm, font=\scriptsize, text=primary!60!black]
     at (-1.0,-3.6)
     {\textbf{Why this displaced recurrence:} (1) no vanishing gradient over distance ---
      every position is one hop away; (2) all positions compute in \emph{parallel}, where
      an RNN must go step by step, which is what made training on internet-scale text
      feasible. The cost is that attention compares every pair, so compute grows with the
      square of sequence length.};
\end{tikzpicture}}%
{Self-attention. Each token's representation is rebuilt as a weighted blend of all tokens,
with the weights learned from content. This mechanism is the basis of every model in
Chapters 20--22.}{attention}

\begin{conceptbox}[Transformer, in Three Sentences]
A transformer is a stack of blocks, each containing self-attention followed by a
small feed-forward network, with residual connections and normalisation to keep
training stable. Because the mechanism has no inherent notion of order, position
information is added explicitly to each token's embedding. Everything in
\chref{genai} and \chref{agentic} is built on this one design.
\end{conceptbox}

\section{Autoencoders}

\begin{definitionbox}[Autoencoder]
A network trained to reproduce its own input through a narrow \term{bottleneck}. It
must therefore learn a compressed \term{latent} representation that keeps only what
matters. Trained only on normal data, it reconstructs normal inputs well and
anomalous inputs badly --- so reconstruction error becomes an anomaly score.
\end{definitionbox}

\dsfig{%
\begin{tikzpicture}[font=\scriptsize,
  l/.style={rounded corners=2pt, draw=#1, fill=#1!12, text=#1!50!black,
            minimum width=0.75cm, font=\tiny\bfseries},
  arr/.style={-{Stealth[length=1.8mm]}, gray!60, line width=0.8pt}]

\node[l=secondary, minimum height=2.4cm] (i) at (0,0)    {};
\node[l=primary,   minimum height=1.7cm] (e1) at (1.5,0) {};
\node[l=primary,   minimum height=1.0cm] (e2) at (3.0,0) {};
\node[l=accent,    minimum height=0.5cm] (z) at (4.5,0)  {};
\node[l=primary,   minimum height=1.0cm] (d1) at (6.0,0) {};
\node[l=primary,   minimum height=1.7cm] (d2) at (7.5,0) {};
\node[l=greenhl,   minimum height=2.4cm] (o) at (9.0,0)  {};

\foreach \a/\b in {i/e1,e1/e2,e2/z,z/d1,d1/d2,d2/o}{\draw[arr] (\a) -- (\b);}

\node[font=\tiny, text=secondary!55!black, align=center] at (0,-1.6)  {input\\$x$};
\node[font=\tiny, text=accent!70!black, align=center]    at (4.5,-1.6){latent $z$\\{\tiny bottleneck}};
\node[font=\tiny, text=greenhl!45!black, align=center]   at (9.0,-1.6){output\\$\hat{x}$};
\node[font=\tiny\bfseries, text=primary!55!black] at (2.25,1.55) {ENCODER};
\node[font=\tiny\bfseries, text=primary!55!black] at (6.75,1.55) {DECODER};

\draw[{Stealth[length=1.8mm]}-{Stealth[length=1.8mm]}, purplehl, line width=1pt]
     (0,-2.15) -- (9.0,-2.15);
\node[font=\tiny\bfseries, text=purplehl] at (4.5,-2.45)
     {loss = reconstruction error $\|x - \hat{x}\|$ --- and \emph{also} the anomaly score};

\node[anchor=west, align=left, text width=3.5cm, font=\tiny, text=gray!55!black] at (10.0,0)
     {Train only on normal data. A reading the network has never seen its like of
      reconstructs badly, and the error flags it.};
\end{tikzpicture}}%
{An autoencoder. The bottleneck forces a compressed representation; the reconstruction
error that trains it also serves, at inference time, as an unsupervised anomaly score ---
which is why this architecture recurs in Chapters 25 and 26.}{autoencoder}

\section{Transfer Learning}

\begin{definitionbox}[Transfer Learning]
Take a model already trained on a large general dataset, keep its learned feature
extractor, and re-train only the final layers on your small specific dataset.
Because early layers learn generic structure (\dsref{cnn-pipeline}), they transfer.
\end{definitionbox}

\begin{worked}[Why Transfer Learning Changes What Is Possible]
You must classify 800 photographs of industrial components as defective or sound.

\smallskip
\textbf{From scratch:} a usable CNN has millions of parameters. With 800 images it
will memorise them within a few epochs. Realistic accuracy: poor, and unstable
between runs.

\textbf{With transfer learning:} take a network pre-trained on ImageNet (1.2 million
images, 1{,}000 categories). Freeze everything but the last block. Replace the
1{,}000-way output with a 2-way one. Now you are fitting a few thousand parameters,
not a few million --- and the frozen layers already detect edges, textures, corners
and surface irregularities, which is most of what ``defective'' means visually.

\smallskip
\textbf{Typical outcome:} 800 images is comfortably enough, training takes minutes
on a laptop rather than days on a GPU cluster, and the result is far more stable.

\textbf{The general rule:} unless you have hundreds of thousands of labelled
examples, do not train a deep vision or language model from scratch. Fine-tune.
This is the single most practical technique in this chapter.
\end{worked}

\rowcolors{2}{white}{lightbg}
\begin{longtable}{@{}p{2.9cm}p{4.0cm}p{3.4cm}p{4.0cm}@{}}
\toprule
\rowcolor{primary!15}
\textbf{Architecture} & \textbf{Assumes} & \textbf{Input} & \textbf{Use for} \\
\midrule
\endhead
Feed-forward (MLP) & Nothing about structure & Fixed-length vectors & Tabular data (but see \chref{features}) \\
CNN & Nearby values are related & Images, spectrograms, sensor grids & Vision, audio, some time series \\
RNN / LSTM / GRU & Order matters; recent past matters most & Sequences & Older sequence work; still fine for short series \\
Transformer & Any element may relate to any other & Sequences, text, images & Language, and increasingly everything else \\
Autoencoder & Normal data lies on a low-dimensional manifold & Anything & Compression, denoising, \textbf{anomaly detection} \\
\bottomrule
\end{longtable}

%---------------------------------------------------------------------
\begin{fourlenses}{Choosing an Architecture}
\lensLA{An \textbf{LSTM or small transformer} over the weekly clickstream
\emph{sequence} captures trajectory --- ``engaged then collapsed'' versus ``steadily
low'' --- which aggregate features flatten away. Worth the complexity only if you
have several cohorts of history; with one cohort, the engineered trend features of
\chref{features} are both better and defensible.}

\lensPM{No deep architecture is appropriate at this data scale. The one legitimate
use is a \textbf{pre-trained language model} applied to unstructured text ---
clustering requirement documents, or classifying retrospective notes by theme.
That is transfer learning doing the work, not your 40 sprints.}

\lensIOT{\textbf{CNNs on spectrograms} are the standard for vibration and acoustic
monitoring: convert the waveform to a time--frequency image and the problem becomes
a vision problem, with all of transfer learning available. \textbf{Autoencoders}
handle the harder case where no failures have ever been recorded --- train on normal
operation and score reconstruction error. Both must then be shrunk to fit an edge
device (\chref{cloudedge}).}

\lensSEC{\textbf{Autoencoders and sequence models} over session logs detect
behaviour that no signature covers. CNNs applied to raw malware bytes treated as an
image are a real and effective technique. But note the adversarial exposure: an
attacker who can query your model can craft inputs that sit just inside the normal
region, and deep models are more susceptible to this than simple ones ---
\chref{cybersecurity} covers the defences.}
\end{fourlenses}

%---------------------------------------------------------------------
\begin{lab}[Transfer Learning and an Autoencoder]
\begin{lstlisting}[language=Python]
# ---------- 1. Transfer learning: 800 images is enough ------------------
from tensorflow import keras

base = keras.applications.MobileNetV2(
    input_shape=(224, 224, 3), include_top=False, weights="imagenet")
base.trainable = False                       # freeze the learned features

model = keras.Sequential([
    base,
    keras.layers.GlobalAveragePooling2D(),
    keras.layers.Dropout(0.3),
    keras.layers.Dense(1, activation="sigmoid"),
])
model.compile(optimizer=keras.optimizers.Adam(1e-3),
              loss="binary_crossentropy", metrics=["AUC"])
# model.fit(train_ds, validation_data=val_ds, epochs=15)

# Then unfreeze the top block and fine-tune at a LOW learning rate:
# base.trainable = True
# for layer in base.layers[:-20]: layer.trainable = False
# model.compile(optimizer=keras.optimizers.Adam(1e-5), ...)   # 1e-5, not 1e-3

# ---------- 2. Autoencoder for unsupervised anomaly detection ----------
import numpy as np
n_feat = X_normal.shape[1]
ae = keras.Sequential([
    keras.layers.Dense(32, activation="relu", input_shape=(n_feat,)),
    keras.layers.Dense(8,  activation="relu"),        # bottleneck
    keras.layers.Dense(32, activation="relu"),
    keras.layers.Dense(n_feat),                       # reconstruct
])
ae.compile(optimizer="adam", loss="mse")
ae.fit(X_normal, X_normal, epochs=50, validation_split=0.1, verbose=0)
#      ^^^^^^^^^ input and target are the same, and only NORMAL data is used

err = np.mean((X_all - ae.predict(X_all, verbose=0))**2, axis=1)
threshold = np.percentile(err[is_normal], 99)      # set on normal data only
print("flagged:", (err > threshold).sum())
\end{lstlisting}
\end{lab}

%---------------------------------------------------------------------
\begin{checkpoint}
\begin{tightnum}
  \item A $3\times3$ convolution kernel is applied to a $32\times32$ image. How many
        weights does the kernel have, and how many would a fully connected layer
        producing the same number of outputs need (roughly)?
  \item Why can attention relate the first and hundredth word of a sentence more
        easily than an LSTM can?
  \item You have 600 labelled images. Argue for transfer learning over training from
        scratch in three sentences.
\end{tightnum}
\end{checkpoint}

%---------------------------------------------------------------------
\begin{chaptersummary}
\begin{tight}
  \item An architecture encodes a structural assumption, which is why the right one
        needs far less data.
  \item Convolution shares one small kernel across all positions, giving locality,
        translation invariance and a huge parameter saving.
  \item RNNs carry a hidden state through time but forget; LSTMs add gates and an
        additive cell state so gradients survive.
  \item Attention lets every position consult every other in one step and computes
        in parallel --- the two reasons transformers displaced recurrence.
  \item Autoencoders compress through a bottleneck; their reconstruction error is a
        ready-made unsupervised anomaly score.
  \item Transfer learning makes deep models viable on small datasets, and should be
        your default for images and text.
\end{tight}
\end{chaptersummary}

\begin{reviewq}
  \item \textbf{(MCQ)} Weight sharing in a CNN means: (a)~layers share a learning
        rate (b)~the same kernel is applied at every position (c)~neurons share
        biases (d)~the model is smaller after pooling.
  \item \textbf{(MCQ)} The main practical advantage of transformers over LSTMs in
        training is: (a)~fewer parameters (b)~parallel computation across positions
        (c)~no need for data (d)~guaranteed convergence.
  \item \textbf{(Short)} Explain the vanishing gradient problem in RNNs and how the
        LSTM cell state addresses it.
  \item \textbf{(Short)} Describe how an autoencoder detects anomalies without ever
        seeing an anomaly during training.
  \item \textbf{(Short)} Why does fine-tuning use a much lower learning rate than
        the initial training of the added layers?
  \item \textbf{(Applied)} You must detect bearing faults from vibration waveforms
        sampled at 20 kHz, with 3{,}000 labelled recordings. Propose an architecture
        and a preprocessing step, and justify both.
  \item \textbf{(Applied)} For each of CNN, LSTM, transformer and autoencoder, name
        one of the four course domains where it is the natural choice, and state the
        structural assumption that makes it so.
\end{reviewq}
